\chapter{System Context}
\ChapterLabel{SystemContext}
\section{Chapter Responsibility}
\label{Section:SystemContextResponsibility}
Define the mission, system boundary, stakeholders, and external result boundary.

\noindent\textbf{Expected result:} An explicit measurement-usability mission, boundary, and stakeholder context.

\index{measurement usability}\index{ADC0804}\index{FT245R}\index{Operator}\index{Integrator}\index{Maintainer}
\section{Mission and Boundary}
\label{Section:MeasurementMission}
The mission is to determine whether a measurement can be used. In
Figure~\ref{Figure:SensorContext}, the system of interest includes the
acquisition hardware and the C++ control app on macOS. The sensor, ADC0804,
clock/counter, interface logic, and FT245R are external to the control app.
The physical quantity enters the system at the sensor. The downstream
consumer is external to the system and receives the published result.
The arrows show the direction and content of the measurement flow. See
\ArchitectureRef{ARC-103} and \ArchitectureRef{ARC-104} for the allocation
and its requirement links.
\begin{figure}[!htbp]
\centering
\begin{tikzpicture}[
  box/.style={draw=s8sBlue,fill=s8sLightBlue,align=center,
    minimum width=2.9cm,minimum height=0.6cm,font=\small},
  boundary/.style={draw=s8sGray,dashed,inner sep=0.25cm},
  flow/.style={->,>=Latex},
  data/.style={font=\footnotesize,align=left,inner sep=3pt}
]
\node[box,draw=s8sGray,fill=white] (environment) at (-2.6,0) {Measured environment};
\node[box] (sensor) at (-2.6,-1.6) {Sensor};
\node[box] (adc) at (-2.6,-2.8) {ADC0804};
\node[box] (counter) at (2.6,-2.8) {Clock/counter};
\node[box,minimum width=8.1cm] (logic) at (0,-4.0) {Interface logic};
\node[box] (ftdi) at (0,-5.4) {FT245R};
\node[box] (host) at (0,-7.0) {Control app\\C++ on macOS};
\node[box,draw=s8sGray,fill=white] (consumer) at (0,-9.0) {Downstream consumer};

\coordinate (right edge) at (5.2,-4);
\node[boundary,fit=(sensor)(adc)(counter)(logic)(ftdi)(host)(right edge)] (system) {};
\node[anchor=north east,font=\footnotesize,align=right] at (system.north east)
  {System of interest};

\draw[flow] (environment.south) -- node[data,right]{Physical quantity} (sensor.north);
\draw[flow] (sensor.south) -- node[data,right]{Analog signal} (adc.north);
\draw[flow] (adc.south) -- node[data,right]{8-bit conversion code}
  (adc.south |- logic.north);
\draw[flow] (counter.south) -- node[data,right]{Counter value}
  (counter.south |- logic.north);
\draw[flow] (logic.south) -- node[data,right]{Associated conversion code and\\counter value via parallel FIFO} (ftdi.north);
\draw[flow] (ftdi.south) -- node[data,right]{Associated conversion code and\\counter value via USB} (host.north);
\draw[flow] (host.south) -- node[data,right,pos=0.66]{Published result\\Code, status, timing} (consumer.north);
\end{tikzpicture}

\caption{Measurement flow across the acquisition hardware, control app, and downstream consumer boundaries.}
\label{Figure:SensorContext}
\end{figure}

\section{Boundary Information and Open Questions}
\label{Section:ContextBoundaryInformation}
\begin{description}[leftmargin=0pt,labelindent=0pt,labelsep=0pt,style=nextline]
\item[Physical quantity to sensor to ADC0804.]
The sensor responds to the physical quantity and supplies an analog signal
to the ADC0804. Identify the sensor and quantity, then establish the
operating conditions, signal units and range, conversion reference, and
calibration needed to interpret the conversion code
(\hyperref[Question:Q-101]{Q-101}, \hyperref[Question:Q-102]{Q-102}).

\item[ADC0804 and clock/counter to interface logic.]
The ADC0804 supplies an 8-bit conversion code \cite{ADC0804Datasheet}.
The clock/counter supplies a counter value through a separate logical input.
Define the counter width, clock, capture event, reset, rollover, and
relationship to measurement time. Establish how the interface logic associates
the code with the correct counter value and how much capture skew is allowed
(\hyperref[Question:Q-111]{Q-111}, \hyperref[Question:Q-112]{Q-112},
\hyperref[Question:Q-104]{Q-104}).

\item[Interface logic to FT245R.]
The interface logic sends the associated conversion code and counter value
through the FT245R parallel FIFO data path. The FT245R datasheet identifies
the D0--D7 data bus \cite{FT245RDatasheet}. Engineering must establish
electrical compatibility, bus ownership, handshake, transfer order,
record framing, and handling of incomplete transfers and backpressure
(\hyperref[Question:Q-102]{Q-102}, \hyperref[Question:Q-103]{Q-103},
\hyperref[Question:Q-112]{Q-112}).

\item[FT245R to control app.]
USB carries the associated conversion code and counter value to the control
app. Define the macOS driver and API, reception and recovery behavior, and
the evidence and decision rules the control app uses to assess measurement
usability. Establish what observations support connection and operating-health
reporting and how the Maintainer receives that information
(\hyperref[Question:Q-106]{Q-106}, \hyperref[Question:Q-107]{Q-107},
\hyperref[Question:Q-113]{Q-113}, \hyperref[Question:Q-115]{Q-115}).

\item[Control app to downstream consumer.]
The control app qualifies measurement usability and publishes code, status,
and timing. Identify the receiving application or system and define the code
meaning, status semantics, timing origin and units, schema and version,
delivery transport, and incomplete or error result policy with the consumer
(\hyperref[Question:Q-114]{Q-114}).
\end{description}

The Operator's decision presentation and any configuration or reverse command
flow remain open in \hyperref[Question:Q-105]{Q-105}.

\section{Stakeholders}
\label{Section:MeasurementStakeholders}
The following needs were supplied by the user. Stakeholder roles are distinct from engineering ownership in \hyperref[Table:ProjectOwners]{the owner register}.

\small
\begin{tabularx}{\linewidth}{L{2.5cm}Y}
\caption{Measurement Stakeholder Needs\label{Table:StakeholderNeeds}}\\
\toprule
Stakeholder & Need and story\\\midrule
\endfirsthead
\toprule
Stakeholder & Need and story\\\midrule
\endhead
\bottomrule
\endlastfoot
\phantomsection\label{Stakeholder:Operator}Operator & Cares about the decision: whether the measurement can be used. \UserStoryRef{US-101}.\\
\phantomsection\label{Stakeholder:Integrator}Integrator & Cares about the meaning and format of the published result. \UserStoryRef{US-102}.\\
\phantomsection\label{Stakeholder:Maintainer}Maintainer & Needs to know whether the source is connected and operating properly. \UserStoryRef{US-103}.\\
\end{tabularx}
\normalsize

\section{Stage and Allocation Boundaries}
This stage remains documentation only; no circuit, firmware, host implementation, or product evidence is produced. macOS is the current host scope, replacing the earlier macOS/Windows/Linux scope. Both software applications retain the earlier C++ direction. The execution target and responsibility of embedded software relative to interface logic remain \hyperref[Question:Q-116]{Q-116}. Selecting ADC0804 and FT245R does not select a processor, FPGA, microcontroller, multiplexer, or host driver. The control app is the host application throughout this book; its name does not imply reverse commands or a graphical UI.
